\chapter{Những quyết định nghiệp vụ đáng chú ý}
\label{ch:quyet-dinh}

Những chương trước mô tả hệ thống làm gì. Chương này nói về sáu quyết định mà chọn một đường
là bỏ một đường khác. Cái giá của mỗi lựa chọn đều đã biết trước, chứ không phát hiện ra lúc
dùng.

Mỗi quyết định có ba phần: bối cảnh, quyết định đã chốt, và cái giá đi kèm. Phần thứ ba đáng
đọc nhất, vì một quyết định không mất gì thường là quyết định chưa gặp tình huống khó.

Riêng mục~\ref{sec:hai-luat-nghich} không theo khuôn ấy. Nó không phải quyết định thứ bảy mà là
chỗ hai quyết định trước gặp nhau, và nhìn qua thì chúng mâu thuẫn.

\section{Duyệt khoá nội dung, nhưng bỏ duyệt được}

\paragraph{Bối cảnh.} Mô tả nghiệp vụ ban đầu coi duyệt và phát hành là một bước, nên một đề chỉ
có hai trạng thái: chưa duyệt và đã phát hành. Thiết kế thực tế tách chúng ra, và cần thêm một
trạng thái mà mô hình cũ không có chỗ chứa: đề vừa tạo ra thì chưa có câu nào.

\paragraph{Quyết định.} Một đề đi qua bốn trạng thái: \emph{trống}, \emph{có câu hỏi},
\emph{đã duyệt}, \emph{đã phát hành}. Duyệt \textbf{khoá nội dung} --- ở trạng thái đã duyệt,
không thêm và không sửa câu hỏi được. Nội dung bị khoá gồm cả lời giải và ánh xạ từ phương án
nhiễu sang lỗi, không chỉ câu hỏi và đáp án. Bỏ duyệt đưa đề về lại nháp, và được phép chừng nào
đề chưa phát hành; cài đặt phát hành đã nhập thì giữ nguyên.

Lý do khoá thì dễ thấy: duyệt là lúc giáo viên nhận trách nhiệm về nội dung. Nếu nội dung còn sửa
được sau khi duyệt thì việc duyệt không có nghĩa gì.

Lý do \emph{đảo ngược được} mới là phần đáng nói. Đây không phải cổng cuối cùng --- phát hành
mới là. Nếu cả hai cổng đều một chiều thì giáo viên gặp hai cửa nặng liên tiếp, và sẽ học cách
bấm qua cả hai cho nhanh. Làm nặng cổng duyệt như vậy là làm hỏng đúng cái cổng thật sự cần sức
nặng.

\paragraph{Cái giá.} Giao diện không được mời giáo viên sửa nữa, khi đề đã duyệt. Nút sửa trên từng
câu phải \textbf{biến mất}, không phải mờ đi: một nút mờ vẫn là một lời mời, và chừng nào nó còn
đó thì câu \emph{nội dung đã khoá} chỉ là lời nói suông. Trạng thái trống cũng phải có một màn
hình rỗng riêng và phải chặn phát hành, nên nó là một trạng thái thật chứ không phải một phép
đếm câu hỏi.

\section{Giờ đóng chặn việc vào, không chặn việc nộp}

\paragraph{Bối cảnh.} Mục~\ref{sec:ba-con-so} đã nói luật này và ví dụ đi kèm. Chỗ đáng bàn ở
đây là vì sao nó được chọn, khi mà nó đi ngược với cách hầu hết người đọc hiểu chữ \emph{đóng}.

\paragraph{Quyết định.} Giờ đóng là hạn \textbf{vào tham gia}, không phải hạn nộp. Mốc này bao
gồm chứ không loại trừ: vào tới hết giờ đóng vẫn được. Ai đã vào rồi thì không bị dừng giữa
chừng.

Mốc chọn kiểu bao gồm, vì một học sinh bấm vào đúng giây cuối không nên bị từ chối bởi một ranh
giới em không nhìn thấy. Còn người đã vào thì không bị dừng, vì ở giai đoạn 1 thời gian làm bài
là thứ hệ thống \emph{cấp cho riêng mỗi em}, đo từ lúc em vào. Cắt ngang là lấy lại một thứ đã
hứa.

\paragraph{Cái giá.} Hệ thống nhận lấy nghĩa vụ giải thích. Vì hiểu nhầm ở đây gây hại theo
chiều ngược với dự đoán, câu giải thích phải xuất hiện ở cả ba nơi --- lúc đang chọn giờ, lúc xác
nhận, và trong biên bản sau khi phát hành --- và phải giống hệt nhau từng chữ. Ba cách diễn đạt
cho một luật sẽ được hiểu thành ba luật.

Kéo theo một ràng buộc nữa: mốc nộp cuối trong câu đó --- 18:15 ở ví dụ
mục~\ref{sec:ba-con-so} --- phải được tính từ giờ đóng cộng thời gian làm bài, chứ không viết
sẵn. Viết cứng thì nó sai ngay lần đầu có người đổi thời gian làm bài, và sai một cách im lặng.

\section{Ba cổng, và ranh giới cứng của cổng đầu vào}

\paragraph{Bối cảnh.} Mục~\ref{sec:ba-cong} đã nêu ba điểm giáo viên bắt buộc tham gia. Hai
trong ba nằm ở đầu ra, tức sau khi trợ lí đã làm xong việc. Cổng còn lại nằm ở đầu vào, và nó là
cổng dễ bị lạm dụng nhất.

\paragraph{Quyết định.} Cổng đầu vào --- trợ lí hỏi lại khi chưa đủ thông tin --- có một ranh giới
cứng:

\begin{itemize}
  \item Chỉ dùng để hỏi thông tin \textbf{trước khi làm một việc đảo ngược được}.
  \item \textbf{Không bao giờ} dùng để xin phép một việc không thu hồi được. Câu hỏi
        \emph{phát hành cho lớp nào} phải đi qua biểu mẫu phát hành và hộp xác nhận.
  \item Trợ lí không đánh dấu lựa chọn nào là nên chọn, khi đó là một quyết định sư phạm. Mỗi
        lựa chọn tự nói ra cái giá của nó.
  \item Luôn giữ lối thoát trả lời tự do. Ba lựa chọn không bao giờ phủ hết.
\end{itemize}

Ranh giới cứng còn quan trọng hơn bản thân việc có cổng. Nếu một bộ ba lựa chọn được phép xin
phép một việc không thu hồi được, thì cổng duyệt bị đi vòng bằng một cú bấm --- và nó \emph{sẽ}
bị đi vòng, vì ba cái nút nhẹ hơn hẳn một biểu mẫu sáu trường.

Không đánh dấu nên chọn, vì chọn nguồn câu hỏi hay mức độ khó là quyết định sư phạm. Trợ lí biết
kho câu hỏi còn gì ở mức nào; quyền chọn thì thuộc về giáo viên. Đưa thông tin là đủ, còn nghiêng
người đọc về một phương án là làm thay việc của họ.

\paragraph{Một cái bẫy đã gặp.} Gợi ý ngầm còn tệ hơn gợi ý công khai. Bộ ba lựa chọn soạn lần
đầu có một phương án duy nhất không kèm mệnh đề nhượng bộ, nên nó nghe nhẹ nhất trong ba phương
án, mà người đọc không nhận ra mình đang được hướng về đó. Vì vậy mọi lựa chọn phải nêu cái giá
bằng cùng một ngữ pháp.

\paragraph{Cái giá.} Cổng đầu vào làm trợ lí chậm lại: nó phải nhận ra mình chưa đủ thông tin
thay vì chọn một mặc định hợp lý. Đổi lại, mỗi lần đoán sai đều tốn một vòng sửa ở đầu ra, nơi
đắt hơn.

Và có một lỗ đã biết, ghi ra đây thay vì để nó âm thầm: \textbf{câu biến thể ở giai đoạn 2 không
đi qua cổng nào}. Đây là ngoại lệ của cổng thứ nhất, tức cổng duyệt. Câu biến thể sinh ra lúc học
sinh đang ngồi làm bài, nên không có khoảng trống nào để chờ một thao tác của giáo viên.

Đó là lý do cơ học chứ không phải lý do nguyên tắc: cổng vẫn đúng, chỉ là chưa có cách giữ nó mà
không bắt cả lớp dừng lại. Vì vậy câu \emph{ba cổng} không còn đọc được như một lời hứa tuyệt
đối, và chỗ nào nhắc tới nó cũng phải trả lời được một câu: thế còn câu biến thể?

\section{Hai đồng hồ ở giai đoạn 2, và chỉ một cái chạy khi học sinh đang học}

\paragraph{Bối cảnh.} Giai đoạn 2 cần thời gian, nhưng thời gian của nó không giống giai đoạn 1 ở
bất cứ điểm nào. Số câu phải làm thay đổi theo từng lượt, phần học sinh đọc lời giải thì không
nên bị hối, và cả giai đoạn có thể kéo dài tới cuối ngày.

\paragraph{Quyết định.} Mục~\ref{sec:hai-dong-ho} đã mô tả cơ chế: số phút mỗi câu là một tỉ lệ
chứ không phải một khoảng, thời lượng một lượt là một quỹ chung, phần giải thích không tính giờ,
và đồng hồ chỉ chạy khi học sinh tự bấm nút làm bài mới.

Ba lý do đứng sau ba lựa chọn ấy. Tỉ lệ thay vì khoảng cố định, vì số câu phải chữa giảm dần qua
từng lượt, và một khoảng cố định sẽ quá rộng ở lượt cuối, quá chật ở lượt đầu. Phần giải thích
không tính giờ, vì đó là lúc học sinh đang học chứ không phải lúc em đang chứng minh; đặt đồng hồ
lên một cuộc giải thích chỉ dạy em gật cho xong rồi bấm tiếp. Ranh giới đặt ở nút làm bài mới, vì
đó là hành động do chính em chọn.

\paragraph{Cái giá.} Ba thứ. Hai thứ đổ lên giáo viên, một thứ làm hệ thống mất khả năng đo.

Giáo viên phải đặt một con số mà họ chưa từng phải nghĩ tới. Không có kinh nghiệm dạy học nào
giúp đoán mấy phút là đủ cho một câu chữa, và đặt sai thì hoặc lớp không kịp, hoặc bài kiểm tra
kéo lê.

Số phút mỗi câu là một con số cho cả đề, nên một câu dài và một câu ngắn được cấp thời gian như
nhau. Muốn phân biệt câu dài câu ngắn thì phải thêm một khái niệm khác, chứ không phải tăng con
số này lên.

Và tổng thời gian một học sinh thực sự bỏ ra \textbf{không dự đoán được}, vì phần giải thích
không có giới hạn. Mọi báo cáo sau này muốn nói em học bao lâu đều phải đo, không suy ra được từ
cài đặt phát hành.

\section{Hai luật nghịch nhau, và vì sao đó là chủ ý}
\label{sec:hai-luat-nghich}

Ở giai đoạn 1, học sinh đã vào thì không bị dừng giữa chừng. Ở giai đoạn 2, hết hạn là dừng, kể
cả khi em đang làm dở một lượt, và câu chưa xong nhận 0 điểm. Hai luật này nghịch nhau, và đó là
chỗ dễ bị một người đọc kỹ coi là thiếu nhất quán rồi đi sửa một trong hai.

Chúng không mâu thuẫn, vì thời gian ở hai giai đoạn là hai thứ khác nhau.

Ở giai đoạn 1, thời gian làm bài là một lượng \textbf{hệ thống cấp cho từng người}, bắt đầu đếm
từ lúc em vào. Cắt ngang một người đang làm là lấy lại một thứ đã hứa với riêng em, và em không
có cách nào biết trước mình sẽ bị lấy mất bao nhiêu.

Ở giai đoạn 2 thì ngược hẳn. Thời lượng một lượt do \textbf{học sinh tự chọn lúc nào bắt đầu},
còn hạn kết thúc là một mốc chung cho cả lớp mà em nhìn thấy suốt cả ngày. Không cắt thì hạn ấy
không tồn tại: ai cũng có thể mở một lượt dài lúc chỉ còn một phút, và kéo bài kiểm tra qua đêm.

Cái giá của việc cắt được trả bằng một nghĩa vụ trên giao diện: màn hình phải \textbf{nói trước}
rằng lượt này có thể bị dừng, ngay tại nút làm bài mới, kèm số phút còn lại thật. Một học sinh bị
dừng giữa chừng mà không được báo trước sẽ tin là hệ thống hỏng.

\section{Lời giải nhiều cách, và mỗi phương án nhiễu gắn một lỗi}

\paragraph{Bối cảnh.} Giai đoạn 2 đòi trợ lí giải thích lỗi rồi dạy lại. Muốn làm được, nó phải
biết hai thứ: học sinh sai vì cái gì, và cách làm đúng là gì. Với một hệ thống chỉ có đáp án
đúng trong tay thì nó không biết cả hai, và chỉ còn đường đoán. Mục~\ref{sec:nhan-loi} đã nói vì
sao đoán sai rồi đi giải thích một lỗi học sinh không mắc thì tệ hơn im lặng.

\paragraph{Quyết định.} Mỗi câu hỏi mang theo lời giải có \textbf{nhiều hơn một cách làm}. Mỗi
phương án nhiễu mang theo \textbf{lỗi mà nó đại diện}, soạn cùng câu hỏi chứ không suy ra lúc
chấm. Và trong một câu, không được có hai phương án viết giống nhau từng chữ.

Luật cuối nghe như một chi tiết nhưng không phải. \emph{Đúng một đáp án đúng} là một luật về cái
nhãn đánh dấu đáp án, nên hai phương án trùng chữ vẫn thoả nó. Nhưng với học sinh thì câu đó có
hai đáp án đúng. Em chọn phương án không được đánh dấu sẽ bị chấm là sai, dù em đã chọn đúng.

Giá trị thật của quyết định này là nó biến chẩn đoán từ \textbf{suy đoán} thành \textbf{tra
cứu}, và làm việc đó ở đúng chỗ: lúc soạn đề, nơi có một người biết môn học đang ngồi. Lời giải
nhiều cách thì cần cho một lý do khác --- một cuộc giải thích chỉ có một cách sẽ không đi đâu
được khi học sinh nói em vẫn chưa hiểu, và lặp lại to hơn không phải là dạy.

\paragraph{Cái giá.} Đây là chỗ hệ thống trả giá đắt nhất trong cả báo cáo, và cái giá nằm
đúng vào một nguyên tắc lõi.

Ánh xạ soạn sẵn gắn mỗi phương án nhiễu với đúng một lỗi: chọn nhiễu B thì hệ thống nói ngay là
lỗi B, nói chắc nịch, kể cả khi em chọn B vì một lý do khác hẳn. Đổi lấy sự chắc chắn ấy, sản
phẩm \textbf{mất khả năng nói ra rằng nó không chắc} đúng ở chỗ nó đang nói chuyện với học sinh.

Đánh đổi này được nhận vì phương án còn lại tệ hơn theo cách đo được. Giữ chẩn đoán bằng suy đoán
nghĩa là mọi câu của mọi em đều rơi vào hàng đợi chờ giáo viên, và giai đoạn 2 không khởi động
được cho ai. Một chẩn đoán do người soạn đề viết ra sai ít hơn một chẩn đoán do máy
đoán --- nhưng nó không phải là không sai.

Ba hệ quả đi kèm, và cả ba đều đã biết trước:

\begin{itemize}
  \item \textbf{Cổng duyệt nặng hơn hẳn.} Duyệt một đề nay là duyệt cả lời giải và ánh xạ nhiễu.
        Mà lời giải khó duyệt hơn câu hỏi: một câu hỏi sai thì đọc là thấy, còn một lời giải sai
        tinh vi thì phải làm thử mới thấy.
  \item \textbf{Lời giải sai bị khuếch đại.} Một câu hỏi sai làm hỏng một câu. Một lời giải sai
        được dạy lại cho mọi em sai câu đó, cùng lúc. Sai sót trong nội dung không còn tỉ lệ với
        số người gặp nó.
  \item \textbf{Câu đúng ở giai đoạn 1 không bao giờ kiểm được lời giải của nó}, vì nó không vào
        giai đoạn 2. Một lời giải sai chỉ lộ ra qua chính những em bị dạy lại bằng nó.
\end{itemize}

\section{Đề có tác giả, và câu từ chối phải giống hệt câu không tìm thấy}

\paragraph{Bối cảnh.} Luật \emph{lớp thuộc về giáo viên} đã được phát biểu từ sớm, nhưng chưa
từng có chỗ để sống: dữ liệu không có trường nào ghi ai là tác giả. Không có chỗ ghi thì không
câu truy vấn nào áp được luật, và không phép kiểm nào hỏi được là luật có đúng hay không.

Chuyện này chưa gây hại khi hệ thống chỉ có màn hình của học sinh, vì ở đó phạm vi dữ liệu đã bị
cắt sẵn theo bài làm của chính em. Màn hình của giáo viên phá vỡ điều đó: mọi thứ bắt đầu từ một
cuộc hội thoại, nên trợ lí phải tự chọn dữ liệu để đọc. Câu \emph{trợ lí chỉ đọc được dữ liệu của giáo viên đang đăng
nhập} phải thành một điều kiện trong truy vấn, không phải một dòng trong lời dặn gửi cho mô hình.

\paragraph{Quyết định.} Mỗi lớp và mỗi đề có đúng một tác giả, và trường đó bắt buộc có. Giáo
viên chỉ đọc và chỉ sửa được lớp và đề của mình.

Và phần đáng bàn: \textbf{lớp không tồn tại và lớp của giáo viên khác nhận cùng một câu trả
lời}. Không có mã lỗi riêng, không có câu từ chối riêng, không có chênh lệch nào để so.

Nếu hai câu trả lời ấy khác nhau thì chúng là một kênh dò: gõ tên lớp cho tới khi thông báo đổi
giọng là biết được lớp nào có thật ở trường. Với một trợ lí trò chuyện bằng tiếng Việt thì việc
dò ấy rẻ nhất: người dò không phải đọc mã lỗi, chỉ cần nhận ra hai câu trả lời nghe khác nhau. Mà
sự khác nhau đó thì chính trợ lí sẽ nói ra, rất nhiệt tình.

Đây cũng là lý do luật này là luật nghiệp vụ chứ không phải một chi tiết kỹ thuật: nó nói hệ
thống được phép để lộ điều gì về những người không phải người đang hỏi.

\paragraph{Cái giá.} Câu từ chối mất đi phần hữu ích nhất của nó. Nó \textbf{không được} nói
\emph{đề này không phải của bạn}, vì chính câu đó đã tiết lộ rằng đề tồn tại. Một giáo viên gõ
nhầm tên lớp sẽ nhận một câu trả lời không giúp họ nhận ra mình gõ sai, và đó là cái giá đã biết.

Một tác giả chứ không phải nhiều, cũng là một lựa chọn có giá: đồng sở hữu một đề là chuyện chưa
ai yêu cầu, và nó làm mọi phép kiểm \emph{đề này của ai} nặng thêm một bậc. Thêm vào sau thì
được; bỏ đi thì không.

\paragraph{Và tên lớp không phải định danh.} Hai lớp cùng tên được phép tồn tại, nên khi giáo
viên nhắc tới một cái tên mà có nhiều lớp mang tên đó, trợ lí phải hỏi lại chứ không được tự
chọn. Đó chính là cổng đầu vào ở mục~\ref{sec:ba-cong} làm việc đúng chỗ của nó.
